\bookchapter{第五十五章：プロジェクト2——天気チェックアプリ（ネットワーク版）}

\chapterintro{家計簿が「内功の鍛錬」なら、天気アプリは「外の世界へ出る」プロジェクトです——インターネットと付き合うことになります。ここで学ぶのは、無料APIの申し込み、JSONの解析、「読み込み中／成功／失敗」の3つの画面状態の処理、そしてAPIキーをきちんと隠す方法です。これをやり遂げれば、どんな「ネット接続アプリ」でも作れる力が手に入ります。ニュースも、為替も、翻訳も、やることはまったく同じです。}

\booksection{無料の天気APIを申し込む}

個人で練習するなら、OpenWeatherMapか、中国国内の「和風天気（QWeather）」がおすすめです。どちらも無料枠があって、登録すればすぐ使えます。OpenWeatherMapを例に見てみましょう。

\begin{enumerate}
  \item アカウントを登録 API keysページでキー（apikey）を1つ発行する。
  \item 無料版で呼べるのはCurrent Weather（現在の天気）のAPIです。api.openweathermap.org/data/2.5/weather? q=beijing\&appid=あなたのキー\&units=metric。
  \item まずブラウザで一度試して、返ってくるJSONがどんな形か見ておきましょう。どんなAPIにつなぐ前にも必ずやる「予習」です。
\end{enumerate}

\textbf{落とし穴注意：キーをコードに直書きしない}\par

APIキーはあなたの「財布のパスワード」みたいなものです。コードに書いてGitHubにプッシュすれば、数時間のうちにクローラーに収集されて勝手に使われます。正しいやり方は、local.properties（このファイルはgitignoreで無視されています）に置いておき、GradleのbuildConfig経由でコードに読み込むことです。「Android API Key 隠す ベストプラクティス」で検索して、その通りに直せば大丈夫です。

\booksection{ネットワーク層：Retrofit＋JSON解析}

まずAPIが返すJSONに合わせて「データモデル」を定義します。対応関係にないフィールドは、すべて無視してかまいません。

\begin{lstlisting}[language=Kotlin]
// JSON: {"weather":[{"description":"晴れ"}],"main":{"temp":26.5,"humidity":40},"name":"北京"}
data class WeatherResponse(
    val weather: List<Weather>,
    val main: Main,
    val name: String
)
data class Weather(val description: String)
data class Main(val temp: Double, val humidity: Int)

interface WeatherApi {
    @GET("data/2.5/weather")
    suspend fun getWeather(
        @Query("q") city: String,
        @Query("appid") key: String,
        @Query("units") metric: String = "metric",
        @Query("lang") lang: String = "zh_cn"
    ): WeatherResponse
}
\end{lstlisting}

GsonはJSONをこれらのdata classへ自動で詰め込みます。フィールド名が一致していれば、魔法は自動的に起こります。

\booksection{3つの画面状態：読み込み中／成功／失敗}

ネット接続アプリの魂は「状態管理」にあります。ネットが遅い、つながらない、都市名を打ち間違えた……どの場合にも、画面はユーザーにきちんと伝える必要があります。sealed class（シールドクラス）で、美しく表現しましょう。

\begin{lstlisting}[language=Kotlin]
sealed class UiState {
    object Loading : UiState()                    // くるくる回る
    data class Success(val data: WeatherResponse) : UiState()  // 成功
    data class Error(val message: String) : UiState()          // 失敗
}

// ViewModel側：
fun search(city: String) {
    viewModelScope.launch {
        _state.value = UiState.Loading
        _state.value = try {
            UiState.Success(api.getWeather(city))
        } catch (e: Exception) {
            UiState.Error("その都市が見つかりません。別の名前で試してみて？")
        }
    }
}

// 画面側：
when (val s = state) {
    is UiState.Loading -> CircularProgressIndicator()   // 回る
    is UiState.Success -> WeatherContent(s.data)        // 表示
    is UiState.Error -> Text(s.message)                 // エラー表示
}
\end{lstlisting}

このLoading/Success/Errorの三段構えは、ネット接続画面の「万能テンプレート」です。まず暗記してしまえば、この先のどんなネット接続アプリでも同じ形で書けます。

\booksection{実戦の画面コード：天気を描き出す}

状態が決まったら、画面はこうなります（要点の抜粋）。

\begin{lstlisting}[language=Kotlin]
@Composable
fun WeatherContent(data: WeatherResponse) {
    Column(
        Modifier.fillMaxSize().padding(24.dp),
        horizontalAlignment = Alignment.CenterHorizontally
    ) {
        Text(data.name, fontSize = 28.sp, fontWeight = FontWeight.Bold)
        Text(
            "${data.main.temp}°C",
            fontSize = 72.sp,
            fontWeight = FontWeight.Light
        )
        Text(
            data.weather.firstOrNull()?.description ?: "不明",
            fontSize = 20.sp,
            color = Color.Gray
        )
        Spacer(Modifier.height(16.dp))
        Row {
            Text(" 湿度 ${data.main.humidity}%")
            Spacer(Modifier.width(24.dp))
            Text(" 体感 ${data.main.temp}°C")
        }
    }
}
\end{lstlisting}

学ぶ価値のある小さな工夫がいくつかあります。firstOrNull() は最初の要素を安全に取り出します（リストが空かもしれないので）。?: “不明” はいざというときの表示を用意しておく——安全に空を扱うという原則は、画面でも同じです。温度は 72.sp の大きい文字で存在感を出し、「細め」のウェイト（Light）で数字を上品に見せる。天気系アプリの定番レシピです。

\booksection{画面と操作の細かい工夫}

\booksection{画面と操作の細かい工夫}

\begin{itemize}
  \item 検索ボックス：上部にTextFieldを置き、都市名を入れて検索（キーボードのアクションを「検索」に設定して、Enterで検索できるように）。
  \item 天気カード：大きい文字の温度（26.5°）、天気の説明（晴れ）、湿度、都市名。
  \item お気に入りの都市：よく見る都市をDataStoreに保存して、次に開いたときに自動で検索——第六章の復習もかねて。
  \item 引っ張って更新：ComposeではPullToRefreshBox（Material 3が提供）を使います。操作の感触が一気に「プロっぽく」なります。
\end{itemize}

\begin{tcbinfo}{ワンポイント：画像で画面を生き生きと}
天気の種類ごとに絵文字かアイコン（）を割り当て、when(description) で対応づけます。3行のコードで、画面のぬくもりが倍になります。サードパーティの天気アイコンライブラリを読み込んでもいいですが、練習の段階なら絵文字で十分です。

\end{tcbinfo}

\begin{tcbpractice}{やってみよう：天気アプリの検収チェックリスト}
(1) 普通に検索して天気が出る。(2) ネットを切っても、アプリはエラーを上品に表示してクラッシュしない（機内モードで試す！）。(3) 存在しない都市を入れても、やさしい表示が出る。(4) 都市を2つお気に入りに入れて、アプリを再起動しても残っている。(5) 画面を回しても状態が消えない。この5つを全部クリアして、このプロジェクトは卒業です。

\end{tcbpractice}

\begin{tcbsummary}{本章のまとめ}
\begin{itemize}
  \item まずブラウザでAPIを試し、JSONを読んでからコードを書く——「予習」で損をすることはない
  \item APIキーはlocal.propertiesに隠す。リポジトリには絶対入れない
  \item Retrofitインターフェース＋data class：フィールドが合えば、JSONは自動で解析される
  \item Loading/Success/Errorの三段構えは、ネット接続画面の万能テンプレート
  \item 切断、空の入力、存在しない都市——エラー処理はネット接続アプリの合格ライン
\end{itemize}

\end{tcbsummary}
